
\section{Mexanik energiyaning saqlanish qonuni}

\begin{enumerate}

\item Tosh yer yuzasidan 10 m/s tezlik bilan vertikal yuqoriga otildi. Qanday balandlikda toshning kinetik energiyasi 5 marta kamayadi? ($g = 10 \, m/s^2$)
\begin{enumerate}[label=\Alph*)]
\item 2
\item 4
\item 6
\item 8
\end{enumerate}

\item Jism yer yuzasidan 20 m/s tezlik bilan vertikal yuqoriga otildi. Qanday balandlikda shu jismning kinetik energiyasi potensial energiyaga teng bo‘ladi? Yer yuzasidagi potensial energiyani nolga teng deb qabul qiling. ($g = 10 \, m/s^2$)
\begin{enumerate}[label=\Alph*)]
\item 10
\item 5
\item 15
\item 20
\end{enumerate}

\item 0,5 kg massali jism yer yuzasidan 10 m balandlikdan 10 m/s tezlik bilan otildi. Jismning yerga tushish paytidagi kinetik energiyasi qanday bo‘ladi (J)? ($g = 10 \, m/s^2$)
\begin{enumerate}[label=\Alph*)]
\item 25
\item 50
\item 100
\item 75
\end{enumerate}

\item Jism yer yuzasidan 10 m balandlikdan gorizontga burchak ostida 20 m/s tezlik bilan otildi. 25 m balandlikda uning tezligi qanchaga teng bo‘ladi (m/s)? ($g = 10 \, m/s^2$)
\begin{enumerate}[label=\Alph*)]
\item 5
\item 10
\item 15
\item 20
\end{enumerate}

\item Jism 30 m balandlikdan 10 m/s tezlik bilan vertikal pastga otildi. Yer sirtidan qanday balandlikda (m) jismning kinetik energiyasi ikki marta oshadi? ($g = 10 \, m/s^2$)
\begin{enumerate}[label=\Alph*)]
\item 25
\item 15
\item 10
\item 20
\end{enumerate}

\item Agar trayektoriyaning eng yuqori nuqtasida jismning kinetik energiyasi potensial energiyaga teng bo‘lsa, u yer yuzasidan gorizontga qanday burchak ostida otilgan? Yer yuzasidagi potensial energiyani nolga teng deb qabul qiling.
\begin{enumerate}[label=\Alph*)]
\item $30^\circ$
\item $60^\circ$
\item $45^\circ$
\item $90^\circ$
\end{enumerate}

\item Kichkina jism silliq gorizontal tekislik bo‘ylab 4 m/s tezlik bilan sirpanadi va tepalikka ko‘tariladi. Jism tekislik sathidan qanday balandlikka (sm da) ko‘tariladi? ($g = 10 \, m/s^2$)
\begin{enumerate}[label=\Alph*)]
\item 40
\item 60
\item 80
\item 100
\end{enumerate}

\item 5 m uzunlikdagi ipga shar osilgan. Shar ip osilgan nuqta joylashgan balandlikkacha og‘ishi uchun unga qanday gorizontal tezlik berish kerak (m/s)? ($g = 10 \, m/s^2$)
\begin{enumerate}[label=\Alph*)]
\item 5
\item 15
\item 20
\item 10
\end{enumerate}

\item 0,4 m uzunlikdagi vaznsiz qattiq sterjenga osilgan sharcha vertikal tekislikda to‘liq aylanishi uchun unga qanday gorizontal minimal tezlik berish kerak (m/s)? ($g = 10 \, m/s^2$)
\begin{enumerate}[label=\Alph*)]
\item 2
\item 4
\item 6
\item 8
\end{enumerate}

\item 80 sm uzunlikdagi yengil sterjen ikki uchiga mahkamlangan 1 kg va 3 kg li yuklar bilan o‘zining o‘rtasidan o‘tgan gorizontal o‘q atrofida erkin aylana oladi. Sterjen gorizontal vaziyatga keltirilib qo‘yib yuboriladi. Sterjen vertikal holatni egallagan paytda yuklarning tezligini toping (m/s). ($g = 10 \, m/s^2$)
\begin{enumerate}[label=\Alph*)]
\item 2
\item 1
\item 3
\item 4
\end{enumerate}

\end{enumerate}



\begin{table}[h!]
\centering
\begin{tabular}{|c|c|c|c|c|c|}
\hline
\textbf{№} & \textbf{Javob} & \textbf{№} & \textbf{Javob} & \textbf{№} & \textbf{Javob} \\ \hline
1 & B  & 5 & A  & 9 & B  \\ \hline
2 & A  & 6 & C  & 10 & A  \\ \hline
3 & D  & 7 & C  & & \\ \hline
4 & B  & 8 & D  & & \\ \hline
\end{tabular}
\end{table}





\section{Mexanik energiyaning ichki energiyaga o‘tishi. Ishqalanish kuchlarining ishi}

\begin{enumerate}

\item 100 g massali shar 2 m balandlikdan po‘lat plita ustiga erkin tushadi va 1 m balandlikkacha sakraydi. Urilish paytida issiqlik ko‘rinishida yo‘qotilgan energiyani aniqlang (J). Havoning qarshiligi inobatga olinmasin. ($g = 10 \, m/s^2$)
\begin{enumerate}[label=\Alph*)]
\item 0,5
\item 1,5
\item 1
\item 2
\end{enumerate}

\item Agar 500 g massali to‘p polga urilib sakragandan so‘ng 50 sm balandlikka ko‘tarilsa va urilish paytida 2 J energiya ajralsa, to‘p boshlang‘ich tezliksiz qanday balandlikdan (sm da) tushgan? ($g = 10 \, m/s^2$)
\begin{enumerate}[label=\Alph*)]
\item 45
\item 135
\item 180
\item 90
\end{enumerate}

\item 0,1 kg massali to‘p poldan 2 m balandlikdan boshlang‘ich tezliksiz qo‘yib yuborildi. Agar to‘pning polga birinchi va ikkinchi urilishlari orasidagi vaqt 1,2 s bo‘lsa, birinchi urilishda ajraladigan issiqlik miqdorini (mJ da) toping. ($g = 10 \, m/s^2$)
\begin{enumerate}[label=\Alph*)]
\item 100
\item 200
\item 300
\item 400
\end{enumerate}

\item 20 m/s tezlik bilan vertikal yuqoriga otilgan 400 g massali to‘p yerga 15 m/s tezlik bilan tushdi. Havoning qarshilik kuchini yengishga sarflangan ishni aniqlang (J).
\begin{enumerate}[label=\Alph*)]
\item 35
\item 15
\item 25
\item 45
\end{enumerate}

\item 500 g massali to‘p yer sirtidan 15 m balandlikdan 10 m/s tezlik bilan vertikal pastga otildi. To‘p yerga 16 m/s tezlik bilan tushdi. Harakat davomida havoning qarshilik kuchi bajaradigan ishning absolyut qiymatini toping (J). ($g = 10 \, m/s^2$)
\begin{enumerate}[label=\Alph*)]
\item 18
\item 36
\item 54
\item 72
\end{enumerate}

\item Agar tosh yerga tushish paytidagi tezligi 18 m/s va havoning qarshilik kuchini yengishga sarflangan ish 38 J bo‘lsa, tosh qanday balandlikdan tushgan (m)? ($m = 1 \, kg, g = 10 \, m/s^2$)
\begin{enumerate}[label=\Alph*)]
\item 10
\item 30
\item 20
\item 40
\end{enumerate}

\item Avtomobil 72 km/h tezlik bilan harakatlanib kelmoqda. Tog‘ etagida dvigatel o‘chirildi. Avtomobil 5 m balandlikka ko‘tarilib to‘xtadi. Boshlang‘ich kinetik energiyaning qanday qismi (foizda) ishqalanishga sarflandi? ($g = 10 \, m/s^2$)
\begin{enumerate}[label=\Alph*)]
\item 75
\item 25
\item 50
\item 100
\end{enumerate}

\item 4,2 kg massali jism gorizontal kuch ta’sirida gorizantal tekislik bo‘ylab sirpanadi. Jismning tekislikka ishqalanish koeffitsiyenti 0,1. 25 m yo‘lda issiqlik ko‘rinishida ajraladigan energiyani aniqlang (J). ($g = 10 \, m/s^2$)
\begin{enumerate}[label=\Alph*)]
\item 55
\item 105
\item 155
\item 205
\end{enumerate}

\item Qiyalik burchagi gorizontga $45^\circ$ bo‘lgan taxta bo‘ylab 5 kg massali brusok ip yordamida 1 m balandlikka tortib chiqariladi. Brusokning taxtaga ishqalanish koeffitsiyenti 0,3. Taxta va brusokni qizdirishga ketadigan energiyani toping (J). ($g = 10 \, m/s^2$)
\begin{enumerate}[label=\Alph*)]
\item 15
\item 30
\item 45
\item 5
\end{enumerate}

\item Chana tog‘dan tushib, gorizontal yuza bo‘ylab 1 m yo‘lni o‘tadi va to‘xtaydi. Agar chana va yo‘l orasidagi ishqalanish koeffitsiyenti 0,2 bo‘lsa, chananing tog‘ etagidagi tezligini aniqlang (m/s). ($g = 10 \, m/s^2$)
\begin{enumerate}[label=\Alph*)]
\item 4
\item 3
\item 2
\item 1
\end{enumerate}

\end{enumerate}



\begin{table}[h!]
\centering
\begin{tabular}{|c|c|c|c|c|c|}
\hline
\textbf{№} & \textbf{Javob} & \textbf{№} & \textbf{Javob} & \textbf{№} & \textbf{Javob} \\ \hline
1 & C & 5 & B & 9 & A \\ \hline
2 & D & 6 & C & 10 & C \\ \hline
3 & B & 7 & A & & \\ \hline
4 & A & 8 & B & & \\ \hline
\end{tabular}
\end{table}




\section{Mexanik energiyaning tashqi kuchlar ta’sirida o`zgarishi}

\begin{enumerate}

\item Deformatsiyalanmagan prujina yordamida vertikal devorga ulangan brusok silliq polda yotadi. Prujinaning o'qi gorizontal, uning bikirligi 50 N/m. Brusokka prujina o'qi bo'ylab yo'nalgan 4 N o'zgarmas kuch ta’sir eta boshlaydi. Prujinaning maksimal deformatsiyasini (cm da) toping.
\begin{enumerate}[label=\Alph*)]
\item 8
\item 16
\item 24
\item 32
\end{enumerate}

\item Deformatsiyalanmagan prujina yordamida vertikal devorga ulangan 100 g massali brusok silliq polda yotadi. Prujinaning o'qi gorizontal, uning bikirligi 250 N/m. Brusokka prujina o'qi bo'ylab yo'nalgan 4 N o'zgarmas kuch ta’sir eta boshlaydi. Brusokning maksimal tezligini (cm/s da) toping.
\begin{enumerate}[label=\Alph*)]
\item 40
\item 60
\item 80
\item 100
\end{enumerate}

\item Deformatsiyalanmagan prujina yordamida vertikal devorga ulangan 250 g massali brusok polda yotadi. Prujinaning o'qi gorizontal, uning bikirligi 100 N/m, ishqalanish koeffitsiyenti 0,4. Brusokka prujina o'qi bo'ylab yo'nalgan 3 N o'zgarmas kuch ta’sir eta boshlaydi. Prujinaning maksimal deformatsiyasini (cm da) toping. ($g = 10 \, m/s^2$)
\begin{enumerate}[label=\Alph*)]
\item 4
\item 2
\item 6
\item 8
\end{enumerate}

\item Deformatsiyalanmagan prujina yordamida vertikal devorga ulangan 250 g massali brusok polda yotadi. Prujinaning o'qi gorizontal, uning bikirligi 100 N/m, ishqalanish koeffitsiyenti 0,4. Brusokka prujina o'qi bo'ylab yo'nalgan 3 N o'zgarmas kuch ta’sir eta boshlaydi. Brusokning maksimal tezligini (cm/s da) toping. ($g = 10 \, m/s^2$)
\begin{enumerate}[label=\Alph*)]
\item 20
\item 30
\item 50
\item 40
\end{enumerate}

\item Shiftga 100 N/m bikirlikdagi prujina yordamida yuk osilgan. Yukka vertikal pastga yo’nalgan 6 N o’zgarmas kuch ta’sir eta boshlaydi. Yukning maksimal ko’chishini (cm da) toping.
\begin{enumerate}[label=\Alph*)]
\item 6
\item 12
\item 18
\item 24
\end{enumerate}

\item Shiftga 50 N/m bikirlikdagi prujina yordamida 0,5 kg massadagi yuk osilgan. Yukka vertikal pastga yo’nalgan 10 N o’zgarmas kuch ta’sir eta boshlaydi. Yukning maksimal tezligini (m/s) toping.
\begin{enumerate}[label=\Alph*)]
\item 1
\item 4
\item 2
\item 3
\end{enumerate}

\item 100 N/m bikirlikdagi prujina yordamida qiya tekislikning uchi bilan ulangan brusok shu silliq qiya tekislikda yotadi. Brusokka tekislik bo’ylab pastga yo’nalgan 12 N o’zgarmas kuch ta’sir eta boshlaydi. Brusokning maksimal ko’chishini (cm da) toping.
\begin{enumerate}[label=\Alph*)]
\item 24
\item 12
\item 36
\item 48
\end{enumerate}

\item Gorizont bilan burchak hosil qiluvchi ($\sin \alpha = 0,6$) qiya tekislikda 1 kg massali brusok yotadi. Brusok qiya tekislikning uchi bilan 64 N/m bikirlikdagi deformatsiyalanmagan prujina yordamida ulangan. Brusokning tekislikka ishqalanish koeffitsiyenti 0,8. Brusokka tekislik bo’ylab yuqoriga yo’nalgan 14 N o’zgarmas kuch ta’sir eta boshlaydi. Brusok to’xtaguncha qanday masofani (cm da) bosib o’tadi? ($g = 10 \, m/s^2$)
\begin{enumerate}[label=\Alph*)]
\item 2,5
\item 7,5
\item 10
\item 5
\end{enumerate}

\end{enumerate}



\begin{table}[h!]
\centering
\begin{tabular}{|c|c|c|c|}
\hline
\textbf{№} & \textbf{Javob} & \textbf{№} & \textbf{Javob} \\ \hline
1 & B & 5 & B \\ \hline
2 & C & 6 & C \\ \hline
3 & A & 7 & A \\ \hline
4 & D & 8 & D \\ \hline
\end{tabular}
\end{table}